\documentclass[letterpaper]{article}
\usepackage{amsmath,amssymb}
\usepackage{lmodern}
\usepackage{iftex}
\ifPDFTeX
  \usepackage[T1]{fontenc}
  \usepackage[utf8]{vietnam}
\else
  \usepackage{fontspec}
\fi
\usepackage{textcomp} % provide euro and other symbols
\usepackage{lipsum}
\usepackage{setspace}
\usepackage{hyperref}
\usepackage{xcolor}
\usepackage{longtable,booktabs,array}
\usepackage{multirow}
\usepackage{calc} % for calculating minipage widths
% Correct order of tables after \paragraph or \subparagraph
\usepackage{etoolbox}
\usepackage{adjustbox}
\usepackage{graphicx}
\makeatletter
\def\maxwidth{\ifdim\Gin@nat@width>\linewidth\linewidth\else\Gin@nat@width\fi}
\def\maxheight{\ifdim\Gin@nat@height>\textheight\textheight\else\Gin@nat@height\fi}
\makeatother
% Scale images if necessary, so that they will not overflow the page
% margins by default, and it is still possible to overwrite the defaults
% using explicit options in \includegraphics[width, height, ...]{}
\setkeys{Gin}{width=\maxwidth,height=\maxheight,keepaspectratio}
\usepackage{tabularx}
\usepackage{tikz}
\usetikzlibrary{arrows.meta,positioning,calc,fit,shapes.geometric,backgrounds}
\usepackage[top=15mm, bottom=15mm, left=30mm, right=30mm]{geometry} % page layout
% add biblatex
\usepackage[backend=bibtex,style=numeric,natbib=true,sorting=nty,maxbibnames=99]{biblatex}
\addbibresource{ref.bib}
% re-define section/ subsection
\usepackage{titlesec}
\providecommand{\tightlist}{%
\setlength{\itemsep}{0pt}\setlength{\parskip}{0pt}}
\setcounter{secnumdepth}{-\maxdimen} % remove section numbering
\ifdefined\DeclareUnicodeCharacter\DeclareUnicodeCharacter{0301}{}\fi
\hypersetup{colorlinks=true,linkcolor=blue,citecolor=blue,urlcolor=blue}

\setlength{\parindent}{1em}
\setlength{\parskip}{0.25em}
\renewcommand{\arraystretch}{1.15}
\setlength{\tabcolsep}{4pt}
\setlength{\LTpre}{6pt}
\setlength{\LTpost}{6pt}
\setlength{\textfloatsep}{10pt plus 2pt minus 2pt}
\setlength{\intextsep}{8pt plus 2pt minus 2pt}
\setlength{\abovecaptionskip}{5pt}
\setlength{\belowcaptionskip}{0pt}

\titleformat{\section}
{\large\bfseries\color{blue}}
{\thesection}{1em}{}
\titleformat{\subsection}
{\large\bfseries\color{blue}}
{\thesection}{1em}{}
\titleformat{\subsubsection}
{\normalsize\color{blue}}
{\thesection}{1em}{}

\author{Trần Ngô Anh Quân}
\date{Tháng 10 năm 2026}

\begin{document}

\begin{table*}[!ht]
\centering
\begin{tabular}{m{6em} c m{7em} m{5em} }
\multirow{3}{3em}{
\begin{minipage}{0.15\textwidth}
\includegraphics{asset/Figure/image1.png}
\end{minipage}}
& & & \\ \cline{3-4}
& \small ĐẠI HỌC QUỐC GIA TP.HCM & \multicolumn{1}{|c|}{\multirow{2}{8em}{\small Ngày nhận hồ sơ}} & \multicolumn{1}{|c|}{}\\
& \small \textbf{TRƯỜNG ĐẠI HỌC CÔNG NGHỆ THÔNG TIN} & \multicolumn{1}{|c|}{} & \multicolumn{1}{|c|}{}\\\cline{3-4}
& & \multicolumn{2}{|c|}{\textit{\small (Do CQ quản lý ghi)}}\\ \cline{3-4}
\end{tabular}
\end{table*}

\begin{quote}
\centering
\textbf{\Large THUYẾT MINH}

\Large ĐỀ TÀI KHOA HỌC VÀ CÔNG NGHỆ CẤP SINH VIÊN 2026
\end{quote}

\section{A. THÔNG TIN CHUNG}

\subsection{A1. Tên đề tài}
\begin{itemize}
\item[-] Tên tiếng Việt (IN HOA): \textbf{NGHIÊN CỨU VÀ TỐI ƯU HÓA TÍNH CỤC BỘ BỘ NHỚ, CÔ LẬP GIAO THỨC TRONG MÁY CHỦ MẠNG ĐA GIAO THỨC HIỆU NĂNG CAO $\mu$WEBZOCKETS BẰNG ZIG}
\item[-] Tên tiếng Anh (IN HOA): \textbf{DATA-ORIENTED MEMORY LOCALITY AND BOUNDED PROTOCOL CONFINEMENT IN $\mu$WEBZOCKETS: A HIGH-CONCURRENCY MULTI-PROTOCOL SERVER ARCHITECTURE IN ZIG}
\end{itemize}

\subsection{A2. Thời gian thực hiện}
\textbf{12} tháng (kể từ khi được duyệt).

\subsection{A3. Tổng kinh phí}
\emph{(Lưu ý tính nhất quán giữa mục này và mục B8. Tổng hợp kinh phí đề nghị cấp)}

\noindent Tổng kinh phí: \textbf{15.000.000} đồng, gồm:
\begin{itemize}
\item Kinh phí từ Trường Đại học Công nghệ Thông tin: \textbf{15.000.000} đồng
\end{itemize}

\subsection{A4. Chủ nhiệm}
\begin{spacing}{1.5}
\indent \indent Họ và tên: Trần Ngô Anh Quân \par
Ngày, tháng, năm sinh: 15/09/2004 \hspace{2cm} Giới tính (Nam/Nữ): Nam \par
Số CMND/CCCD: 079204012345; Ngày cấp: 15/10/2021; Nơi cấp: Cục CSQLHC về TTXH \par
Mã số sinh viên: 25521506 \par
Số điện thoại liên lạc: 0901234567 \par
Đơn vị (Khoa): Mạng máy tính và Truyền thông dữ liệu \par
Số tài khoản: 255215068888 \hspace{4cm} Ngân hàng: Vietcombank - Chi nhánh Đông Sài Gòn
\end{spacing}

\subsection{A5. Thành viên đề tài}
\begin{longtable}[c]{@{} >{\raggedright\arraybackslash}p{(\columnwidth - 6\tabcolsep) * \real{0.0553}} >{\raggedright\arraybackslash}p{(\columnwidth - 6\tabcolsep) * \real{0.3235}} >{\raggedright\arraybackslash}p{(\columnwidth - 6\tabcolsep) * \real{0.3637}} >{\raggedright\arraybackslash}p{(\columnwidth - 6\tabcolsep) * \real{0.2575}}@{}}
\toprule()
\begin{minipage}[b]{\linewidth}\raggedright \textbf{TT} \end{minipage} &
\begin{minipage}[b]{\linewidth}\raggedright \textbf{Họ tên} \end{minipage} &
\begin{minipage}[b]{\linewidth}\raggedright \textbf{MSSV} \end{minipage} &
\begin{minipage}[b]{\linewidth}\raggedright \textbf{Khoa} \end{minipage} \\
\midrule()
\endhead
1 & Trần Ngô Anh Quân & 25521506 & Mạng máy tính và Truyền thông dữ liệu \\
\bottomrule()
\end{longtable}

\section{B. MÔ TẢ NGHIÊN CỨU}

\indent \indent Nghiên cứu này tập trung giải quyết bài toán cốt lõi trong kỹ thuật hạ tầng mạng hiện đại: rào cản ``Bức tường bộ nhớ'' (\emph{Memory Wall})~\cite{drepper_memory} và hiện tượng bùng nổ độ trễ đuôi (\emph{tail latency}) trong các máy chủ mạng bất đồng bộ khi đối mặt với hàng chục ngàn kết nối đồng thời. Thay vì mô hình hóa kết nối mạng bằng đồ thị đối tượng phân tán trên heap theo lối Hướng đối tượng truyền thống, đề tài tiếp cận theo triết lý Thiết kế Định hướng Dữ liệu (\emph{Data-Oriented Design} -- DoD)~\cite{kelly_dod}, ứng dụng ngôn ngữ lập trình hệ thống thế hệ mới Zig 0.16.0 để xây dựng thư viện máy chủ mạng đa giao thức hiệu năng cao mang tên \textbf{$\mu$WebZockets}. Dự án tích hợp đầy đủ ngăn xếp giao thức chuẩn hóa IETF bao gồm HTTP/1.1~\cite{rfc9112}, HTTP/2~\cite{rfc9113} kèm HPACK~\cite{rfc7541}, WebSocket~\cite{rfc6455} kèm Per-Message Deflate~\cite{rfc7692}, HTTP/3~\cite{rfc9114} và QUIC~\cite{rfc9000} thông qua lsquic~\cite{lsquic}, HTTP QUERY~\cite{rfc10008}, và bảo mật TLS 1.3 qua BoringSSL~\cite{boringssl}.

Kiến trúc đề xuất đã đạt được những minh chứng thực nghiệm ấn tượng trên các đấu trường benchmark quốc tế: xếp hạng hàng đầu với hơn $170{,}000$ req/s tại bộ chuẩn \textit{Zig Web Servers Benchmark}~\cite{bench_zig_bd}, vượt qua 15 bộ kiểm thử tải nặng khắt khe trong \textit{HttpArena}~\cite{httparena_pr1536} với đánh giá hoàn thiện 4/4 tính năng, và vượt qua 514/517 ca kiểm thử chuẩn tuân thủ giao thức WebSocket trong \textit{Autobahn Testsuite}~\cite{autobahn}. Tổng quan cấu trúc kiến trúc đa tầng của hệ thống được minh họa trực quan tại Hình~\ref{fig:overview-arch}.

\begin{figure}[!hbt]
\centering
\begin{tikzpicture}[
    scale=0.92, transform shape,
    box/.style={draw=blue!80, fill=blue!5, thick, rounded corners=3pt, align=center, font=\small},
    workerbox/.style={draw=teal!80, fill=teal!6, thick, rounded corners=4pt, inner sep=6pt},
    slabbox/.style={draw=orange!80, fill=orange!8, thick, dashed, rounded corners=2pt, align=center, font=\footnotesize},
    kbox/.style={draw=purple!80, fill=purple!5, thick, rounded corners=3pt, align=center, font=\small},
    arrow/.style={-{Stealth[length=2.5mm]}, thick, draw=blue!70}
]

% Network Ingress
\node[kbox, text width=12cm, fill=gray!10, draw=gray!70] (network) {
    \textbf{Tầng Mạng Vật Lý \& Hệ Điều Hành (Linux Kernel 6.x)}\\
    \footnotesize Hàng đợi gói tin RX/TX $\bullet$ Phân luồng phần cứng NIC (RSS) $\bullet$ TCP/UDP Sockets (\texttt{SO\_REUSEPORT})
};

% Multi-Worker Shared-Nothing Layer
\node[workerbox, text width=5.6cm, below=0.7cm of network.south west, anchor=north west] (w0) {
    \centering \textbf{\color{teal!90!black}Worker 0 (Ghim CPU Core 0)}\\[3pt]
    \begin{tikzpicture}[node distance=0.35cm, font=\footnotesize]
        \node[box, text width=5cm, fill=teal!10] (loop0) {\textbf{Vòng lặp sự kiện libxev}\\(\texttt{io\_uring} / Completion Ring)};
        \node[slabbox, text width=5cm, below=of loop0] (slab0) {\textbf{Connection Slab tĩnh (Zero-Alloc)}\\$\bullet$ Khe kết nối 64B aligned\\$\bullet$ Freelist thời gian $O(1)$\\$\bullet$ Write Queue Ring Buffer cố định};
        \node[box, text width=5cm, below=of slab0, fill=blue!10] (proto0) {\textbf{Bộ thẩm định giao thức Bounded}\\$\bullet$ HTTP/1.1 \& RFC 10008 QUERY\\$\bullet$ HTTP/2 (SoA Streams, HPACK)\\$\bullet$ WebSocket (SIMD Mask, Deflate)\\$\bullet$ HTTP/3 / QUIC (lsquic C-FFI)};
    \end{tikzpicture}
};

\node[workerbox, text width=5.6cm, below=0.7cm of network.south east, anchor=north east] (w1) {
    \centering \textbf{\color{teal!90!black}Worker $N-1$ (Ghim CPU Core $N-1$)}\\[3pt]
    \begin{tikzpicture}[node distance=0.35cm, font=\footnotesize]
        \node[box, text width=5cm, fill=teal!10] (loop1) {\textbf{Vòng lặp sự kiện libxev}\\(\texttt{io\_uring} / Completion Ring)};
        \node[slabbox, text width=5cm, below=of loop1] (slab1) {\textbf{Connection Slab tĩnh (Zero-Alloc)}\\$\bullet$ Khe kết nối 64B aligned\\$\bullet$ Freelist thời gian $O(1)$\\$\bullet$ Write Queue Ring Buffer cố định};
        \node[box, text width=5cm, below=of slab1, fill=blue!10] (proto1) {\textbf{Bộ thẩm định giao thức Bounded}\\$\bullet$ HTTP/1.1 \& RFC 10008 QUERY\\$\bullet$ HTTP/2 (SoA Streams, HPACK)\\$\bullet$ WebSocket (SIMD Mask, Deflate)\\$\bullet$ HTTP/3 / QUIC (lsquic C-FFI)};
    \end{tikzpicture}
};

% User Application Layer
\node[box, text width=12cm, fill=blue!12, draw=blue!90, below=0.6cm of w0.south west, anchor=north west] (app) {
    \textbf{Tầng Ứng Dụng \& Xử Lý Nghiệp Vụ (Typed Request / Response API)}\\
    \footnotesize Radix-Tree Router không cấp phát $\bullet$ Middleware Pipeline $\bullet$ Zero-Copy Borrowed Memory Slices
};

% Arrows
\draw[arrow] (network.south -| w0.north) -- (w0.north) node[midway, right, font=\tiny] {Gói tin socket};
\draw[arrow] (network.south -| w1.north) -- (w1.north) node[midway, right, font=\tiny] {Gói tin socket};
\draw[arrow] (w0.south) -- (app.north -| w0.south) node[midway, right, font=\tiny] {Dispatch sự kiện};
\draw[arrow] (w1.south) -- (app.north -| w1.south) node[midway, right, font=\tiny] {Dispatch sự kiện};

\end{tikzpicture}
\caption{Kiến trúc tổng thể của máy chủ mạng đa giao thức $\mu$WebZockets: Mô hình cụm Worker không chia sẻ (\emph{Shared-Nothing}), ghim lõi CPU độc lập (\emph{CPU Pinning}), vòng lặp sự kiện hướng hoàn tất (\emph{libxev}), và phân vùng bộ nhớ đệm tĩnh (\emph{Bounded Slab}) loại bỏ hoàn toàn cấp phát động trên đường đi nóng.}
\label{fig:overview-arch}
\end{figure}

\subsection{B1. Giới thiệu về đề tài}

\paragraph{1. Tính cấp thiết của đề tài và bối cảnh công nghệ:}
Trong kỷ nguyên của các hệ thống phân tán thời gian thực, dịch vụ tài chính tần suất cao, truyền phát đa phương tiện tương tác và kiến trúc microservices mật độ cao, máy chủ mạng đóng vai trò là xương sống quyết định thông lượng và độ tin cậy của toàn bộ hệ thống. Với sự tiến hóa vượt bậc của hạ tầng phần cứng, các bộ vi xử lý máy chủ hiện đại sở hữu hàng chục đến hàng trăm lõi tính toán với chu kỳ xung nhịp sub-nanosecond, cùng băng thông mạng đạt ngưỡng $100\text{ Gbps}$ đến $400\text{ Gbps}$. Tuy nhiên, hiệu năng thực tế của các ứng dụng mạng phần lớn lại không được khai thác tương xứng do rào cản cố hữu mang tên \textbf{``Bức tường bộ nhớ'' (\emph{Memory Wall})}~\cite{drepper_memory}. Trong khi chu kỳ xử lý của CPU chỉ mất khoảng $0.3\text{ ns}$, một lần nạp dữ liệu thất bại từ bộ nhớ chính DRAM (\emph{Cache Miss}) phải trả giá bằng một khoảng thời gian chờ trễ lên tới $50\text{--}100\text{ ns}$ (tương đương với hàng trăm chu kỳ CPU bị lãng phí hoàn toàn).

Phần lớn các framework máy chủ mạng phổ biến hiện nay (như Node.js/Express, Python ASGI, Go net/http, hay Java Netty) được thiết kế theo tư duy Lập trình Hướng đối tượng (\emph{Object-Oriented Programming} -- OOP). Trong mô hình này, mỗi kết nối TCP/UDP được khởi tạo thành một tập hợp các thực thể đối tượng phân tán trên heap thông qua các lệnh gọi cấp phát động (\texttt{malloc}/\texttt{new}). Khi một gói tin mạng đến, CPU buộc phải thực hiện duyệt con trỏ (\emph{pointer chasing}) qua nhiều cấu trúc dữ liệu rời rạc (đối tượng kết nối $\to$ bộ đệm luồng $\to$ bảng băm tiêu đề $\to$ máy trạng thái phân tích cú pháp). Lối thiết kế này dẫn đến ba hệ lụy nghiêm trọng:
\begin{enumerate}
    \item \textbf{Phân mảnh bộ nhớ và suy giảm tính cục bộ cache L1/L2:} Các vùng nhớ bị phân tán ngẫu nhiên trong không gian ảo, phá vỡ khả năng nạp trước tự động (\emph{hardware prefetching}) của phần cứng, làm tăng vọt tỉ lệ trượt cache ($M_{L1}$, $M_{L2}$) trên mỗi chu trình xử lý sự kiện I/O.
    \item \textbf{Hiện tượng khóa tranh chấp bộ cấp phát (\emph{Allocator Contention}):} Ở mức hàng chục ngàn kết nối đồng thời, việc hàng chục luồng liên tục gọi cấp phát và giải phóng bộ nhớ trên một heap chia sẻ gây xung đột dữ liệu và hiện tượng tráo dòng nhớ giả (\emph{False Sharing}).
    \item \textbf{Hiện tượng bùng nổ độ trễ đuôi (\emph{Tail Latency}):} Tác động cộng gộp của việc phân mảnh bộ nhớ và các đợt tạm dừng thu gom rác (\emph{Garbage Collection pauses}) trong các ngôn ngữ có runtime tạo ra sự biến thiên độ trễ rất lớn ở các phân vị cao (P99, P99.9), gây suy giảm nghiêm trọng trải nghiệm người dùng cuối.
\end{enumerate}

\paragraph{2. Cách tiếp cận mới của đề tài $\mu$WebZockets:}
Để giải quyết tận gốc rễ các vấn đề trên, đề tài này đề xuất và hiện thực hóa một kiến trúc máy chủ mạng hoàn toàn mới mang tên \textbf{$\mu$WebZockets}, sử dụng ngôn ngữ lập trình hệ thống \textbf{Zig 0.16.0}. Kiến trúc được xây dựng dựa trên bốn trụ cột thiết kế khoa học:
\begin{enumerate}
    \item \textbf{Thiết kế Định hướng Dữ liệu (\emph{Data-Oriented Design} -- DoD) và Zero-Allocation Slabs:} Khởi tạo toàn bộ tài nguyên (khe kết nối, hàng đợi ghi vòng tròn, bộ nhớ phân tích cú pháp) thành các khối liên tục (\emph{contiguous slabs}) căn chỉnh đúng đường biên cache line 64-byte ngay từ giai đoạn khởi động máy chủ. Trên đường đi nóng (\emph{hot path}), hệ thống hoàn toàn không gọi bộ cấp phát động của hệ điều hành, vận hành danh sách rỗi (\emph{Freelist}) với thời gian hằng số $O(1)$.
    \item \textbf{Chuyển dịch cấu trúc luồng HTTP/2 từ AoS sang SoA:} Thay vì sử dụng Mảng các Cấu trúc (\emph{Array of Structures} -- AoS) truyền thống, trạng thái luồng đa hợp HTTP/2 được tổ chức thành Cấu trúc các Mảng (\emph{Structure of Arrays} -- SoA), giúp các trường dữ liệu nóng (trạng thái mở/đóng, stream ID, hạn mức điều khiển luồng) nằm sát nhau trên một dòng nhớ cache L1d, tối đa hóa mật độ dữ liệu hữu ích.
    \item \textbf{Cô lập giao thức có giới hạn (\emph{Bounded Protocol Confinement}):} Toàn bộ các bộ phân tích cú pháp HTTP/1.1, HTTP/2, HPACK, WebSocket, HTTP/3 và WebTransport đều thực thi dưới các ràng buộc kích thước biên cứng (\emph{bounded buffers}), ứng dụng kỹ thuật mượn lát cắt không sao chép (\emph{zero-copy borrowed slices}). Mọi dữ liệu không hợp lệ hoặc vượt ngưỡng dung lượng quy định đều bị từ chối dứt khoát với mã lỗi HTTP chuẩn xác, ngăn chặn triệt để các dạng tấn công cạn kiệt bộ nhớ như Slowloris hay HTTP/2 CONTINUATION flood.
    \item \textbf{Tăng tốc vi lệnh SIMD và Kiến trúc Cụm không chia sẻ (\emph{Multi-Worker Shared-Nothing}):} Ứng dụng các chỉ thị SIMD (SWAR, AVX2, ARM NEON) quét song song $16\text{--}32\text{ bytes}$ mỗi chu kỳ xung nhịp để phát hiện dấu phân cách CRLF và che mặt nạ WebSocket. Mỗi lõi CPU vật lý vận hành một tiến trình worker độc lập ghim cố định (\emph{core pinning}), sở hữu một vòng lặp sự kiện \texttt{libxev} (dựa trên \texttt{io\_uring}) riêng biệt, triệt tiêu hoàn toàn nhu cầu khóa đồng bộ hóa giữa các luồng.
\end{enumerate}

\paragraph{3. Tình hình nghiên cứu trong, ngoài nước và vị thế thực nghiệm:}
\begin{itemize}
    \item \textit{Trên thế giới:} Các nghiên cứu tiên phong như uWebSockets~\cite{uwebsockets} (C++) hay Actix (Rust) đã chứng minh tiềm năng vượt trội của mô hình zero-allocation. Đồng thời, sự ra đời của giao diện bất đồng bộ nhân Linux \texttt{io\_uring}~\cite{io_uring} do Jens Axboe phát triển và thư viện \texttt{libxev}~\cite{libxev} do Mitchell Hashimoto khởi xướng đã mở ra kỷ nguyên I/O dựa trên cơ chế thông báo hoàn tất (\emph{completion-based I/O}) thay thế cho cơ chế sẵn sàng (\emph{readiness-based epoll}).
    \item \textit{Tại Việt Nam và trường ĐH Công nghệ Thông tin (UIT):} Các đề tài nghiên cứu cấp sinh viên về mạng máy tính phần lớn tập trung vào tầng ứng dụng Web hoặc triển khai dịch vụ đám mây cấp cao; rất hiếm các công trình đi sâu vào kiến trúc tầng thấp của hệ điều hành, tối ưu hóa bộ nhớ đệm vi xử lý, kỹ thuật vector hóa SIMD, và đo kiểm độ trễ hạt mịn bằng công nghệ eBPF.
    \item \textit{Cơ sở thực nghiệm và thành tựu đã được kiểm chứng của dự án:}
    \begin{enumerate}
        \item \textbf{Bảng xếp hạng Zig Web Servers Benchmark (\texttt{bench.zig.bd})~\cite{bench_zig_bd}:} Dự án $\mu$WebZockets được đánh giá độc lập cùng với các framework Zig danh tiếng (như \texttt{dusty+zio}, \texttt{httpz}, \texttt{zap}, \texttt{zzz}, \texttt{std.http}). Kết quả đo kiểm ghi nhận $\mu$WebZockets đạt thông lượng vượt trội xấp xỉ $170{,}011\text{ req/s}$ ở kịch bản Plaintext và $158{,}420\text{ req/s}$ ở kịch bản JSON, duy trì độ trễ trung bình cực thấp chỉ $3.46\text{ ms}$, phân vị P99 ở mức $6.57\text{ ms}$ và tỷ lệ thành công tuyệt đối $100\%$.
        \item \textbf{Tích hợp chính thức vào bộ chuẩn HttpArena (Pull Request \#1536)~\cite{httparena_pr1536}:} $\mu$WebZockets đã được thẩm định và phê duyệt tích hợp vào hệ thống benchmark máy chủ HTTP/WebSocket quốc tế HttpArena với điểm số hoàn thiện tính năng tối đa $4/4$ (Radix router, Middleware pipeline, Typed Request/Response API, Multi-protocol support), vượt qua toàn bộ 15 kịch bản tải nặng bao gồm: \texttt{baseline}, \texttt{pipelined}, \texttt{limited-conn}, \texttt{json-comp}, \texttt{json-tls}, \texttt{baseline-h2}, \texttt{baseline-h2c}, \texttt{baseline-h3}, \texttt{echo-ws}, \texttt{echo-ws-pipeline}, \texttt{latency-1m}, \texttt{latency-500k-8cpu}.
        \item \textbf{Chứng nhận tuân thủ chuẩn WebSocket quốc tế Autobahn Testsuite~\cite{autobahn}:} Vượt qua $514$ ca kiểm thử đạt kết quả \texttt{OK} và $3$ ca \texttt{INFORMATIONAL} trên tổng số $517$ ca kiểm thử khắt khe của bộ chuẩn.
    \end{enumerate}
\end{itemize}

\subsection{B2. Mục tiêu, nội dung, kế hoạch nghiên cứu}

\subsection{B2.1 Mục tiêu}
\paragraph{Mục tiêu tổng quát:}
Nghiên cứu, mô hình hóa, tối ưu hóa và đánh giá thực nghiệm toàn diện kiến trúc máy chủ mạng đa giao thức hiệu năng cao $\mu$WebZockets viết bằng ngôn ngữ Zig 0.16.0; chứng minh một cách định lượng vai trò quyết định của Thiết kế Định hướng Dữ liệu (DoD), cấp phát tĩnh trước (Slab Allocation), tổ chức Structure of Arrays (SoA) và vi lệnh SIMD trong việc tối đa hóa tính cục bộ bộ nhớ cache L1/L2 và triệt tiêu độ trễ đuôi ở mức tải đồng thời cực lớn.

\paragraph{Mục tiêu cụ thể:}
\begin{enumerate}
    \item \textbf{Mục tiêu 1:} Phân tích toán học và thiết kế cấu trúc dữ liệu Connection Slab, danh sách rỗi Freelist $O(1)$, và hàng đợi ghi vòng tròn không khóa, đảm bảo loại bỏ hoàn toàn các thao tác cấp phát bộ nhớ động trên chu trình xử lý yêu cầu/phản hồi.
    \item \textbf{Mục tiêu 2:} Xây dựng mô hình tổ chức bộ nhớ SoA cho quản lý luồng HTTP/2 và HPACK dynamic table, so sánh định lượng hiệu quả nạp cache line 64-byte so với mô hình AoS truyền thống.
    \item \textbf{Mục tiêu 3:} Hoàn thiện bộ cơ chế cô lập giao thức có giới hạn (\emph{Bounded Protocol Confinement}) hỗ trợ đầy đủ các chuẩn IETF: HTTP/1.1 (RFC 9112), HTTP/2 (RFC 9113), HPACK (RFC 7541), WebSocket (RFC 6455), Per-Message Deflate (RFC 7692), HTTP/3 (RFC 9114), QUIC (RFC 9000), HTTP QUERY (RFC 10008) và BoringSSL TLS 1.3 (RFC 8446).
    \item \textbf{Mục tiêu 4:} Xây dựng phương pháp đo kiểm hạt mịn tái lập được bằng Linux eBPF và Hardware Performance Counters, đối sánh số liệu thực nghiệm với các bộ chuẩn quốc tế \url{https://bench.zig.bd/} và HttpArena PR \#1536.
\end{enumerate}

\subsection{B2.2 Nội dung và phương pháp nghiên cứu}

\subsubsection{Nội dung 1: Khảo sát và thiết kế kiến trúc bộ nhớ định hướng dữ liệu (Data-Oriented Design) và Zero-Allocation Slabs.}

\paragraph{Phương pháp thực hiện:}
\begin{itemize}
    \item \textit{Phân tích vi kiến trúc phần cứng:} Nghiên cứu cách thức phân cấp bộ nhớ cache L1d ($32\text{--}48\text{ KB}$, độ trễ 4--5 chu kỳ), L2 ($512\text{ KB}\text{--}1\text{ MB}$, độ trễ 12--14 chu kỳ) và kích thước dòng nhớ chuẩn $64\text{ bytes}$ trên các kiến trúc hiện đại x86\_64 và ARM64~\cite{drepper_memory}.
    \item \textit{Thiết kế Connection Slab liên tục:} Thay vì cấp phát từng đối tượng kết nối qua heap, toàn bộ dung lượng cho $N$ kết nối đồng thời được cấp phát tĩnh trong một mảng phẳng liên tục duy nhất ngay khi khởi động. Các trường dữ liệu nóng (\emph{Hot Fields}: file descriptor, trạng thái phân tích cú pháp, con trỏ vòng đệm) được gộp chung trong 64 bytes đầu tiên để vừa khít trong đúng một dòng cache line đơn lẻ. Các trường dữ liệu lạnh (\emph{Cold Fields}: thông tin chứng chỉ SSL, cấu hình timeouts, thống kê tổng) được chuyển sang vùng đệm riêng biệt để không chiếm dụng cache L1d.
    \item \textit{Cấu trúc Structure of Arrays (SoA) cho HTTP/2:} Trong kết nối HTTP/2 đa hợp (\emph{multiplexing}), một kết nối có thể mở đến 100 luồng stream đồng thời. Nghiên cứu thực hiện chuyển đổi cấu trúc quản lý luồng từ AoS (Mảng các struct chứa trạng thái, bộ đệm, cửa sổ điều khiển luồng) sang SoA (các mảng tách biệt song song: \texttt{states: [100]StreamState}, \texttt{window\_sizes: [100]u32}, \texttt{weights: [100]u8}). Khi quét tìm luồng sẵn sàng ghi dữ liệu, CPU chỉ cần nạp mảng \texttt{window\_sizes}, giúp tăng mật độ phần tử trên mỗi cache line lên gấp 8 lần và kích hoạt bộ nạp trước phần cứng (\emph{hardware stream prefetcher}) hoạt động tối ưu (xem minh họa so sánh tại Hình~\ref{fig:memory-layout}).
    \item \textit{Hàng đợi ghi vòng tròn cố định (\emph{Fixed Ring Buffer}):} Mỗi kết nối được gắn cố định một bộ đệm ghi dạng vòng với con trỏ \texttt{head} và \texttt{tail}. Khi ứng dụng gọi \texttt{res.send()}, dữ liệu được ghi trực tiếp vào ring buffer mà không cần cấp phát thêm nút danh sách liên kết.
\end{itemize}

\begin{figure}[!hbt]
\centering
\begin{tikzpicture}[
    scale=0.88, transform shape,
    cacheline/.style={draw=black!70, fill=yellow!12, thick, inner sep=4pt, rounded corners=2pt},
    bytecell/.style={draw=blue!80, fill=blue!10, minimum width=1.1cm, minimum height=0.65cm, align=center, font=\scriptsize},
    coldcell/.style={draw=gray!60, fill=gray!15, minimum width=1.1cm, minimum height=0.65cm, align=center, font=\scriptsize},
    labelstyle/.style={font=\small\bfseries, align=center}
]

% Part A: Traditional OOP / AoS
\node[labelstyle, text=red!80!black] at (3.5, 3.8) {(a) Mô hình Hướng đối tượng OOP truyền thống (Heap Allocations \& AoS)};
\node[draw=red!70, fill=red!5, thick, rounded corners=4pt, fit={(0,0) (7,3.3)}] (oophull) {};

\node[bytecell, fill=red!15] (obj1) at (1.6, 2.6) {Conn Obj A\\(Heap)};
\node[bytecell, fill=red!15] (obj2) at (5.4, 2.6) {Conn Obj B\\(Heap)};
\node[coldcell] (buf1) at (1.0, 1.2) {Buffer A\\(Mismatched)};
\node[coldcell] (hdr1) at (2.4, 0.5) {Headers\\(Pointer)};
\node[coldcell] (buf2) at (4.8, 1.0) {Buffer B\\(Separate)};
\node[coldcell] (hdr2) at (6.0, 0.4) {Headers\\(Fragmented)};

\draw[-{Stealth}, thick, dashed, red!80] (obj1) -- (buf1) node[midway, left, font=\tiny] {Trỏ heap};
\draw[-{Stealth}, thick, dashed, red!80] (obj1) -- (hdr1);
\draw[-{Stealth}, thick, dashed, red!80] (obj2) -- (buf2) node[midway, right, font=\tiny] {Pointer chase};
\draw[-{Stealth}, thick, dashed, red!80] (obj2) -- (hdr2);

\node[font=\scriptsize, align=center, text=red!80!black] at (3.5, -0.4) {Dữ liệu phân mảnh ngẫu nhiên $\to$ Gây trượt cache L1/L2 liên tục};

% Part B: DoD Contiguous Slab & SoA
\begin{scope}[xshift=8.2cm]
\node[labelstyle, text=teal!90!black] at (3.5, 3.8) {(b) Mô hình Data-Oriented Design $\mu$WebZockets (Slab \& SoA)};
\node[draw=teal!70, fill=teal!5, thick, rounded corners=4pt, fit={(0,0) (7,3.3)}] (dodhull) {};

% Cache line 1 (64 bytes)
\node[draw=teal!90, fill=teal!15, thick, minimum width=6.4cm, minimum height=0.8cm, font=\scriptsize\bfseries, align=center] (cl1) at (3.5, 2.5) {
    Dòng Cache L1d thứ 1 (Căn chỉnh 64-byte boundary)\\[2pt]
    \normalfont\tiny \color{blue!80!black} [fd: 4B] $\bullet$ [State: 4B] $\bullet$ [Ring Head/Tail: 8B] $\bullet$ [Hot Flags: 4B] $\bullet$ [Stream Ptrs: 44B]
};

% Cache line 2
\node[draw=teal!90, fill=teal!15, thick, minimum width=6.4cm, minimum height=0.8cm, font=\scriptsize\bfseries, align=center] (cl2) at (3.5, 1.3) {
    Dòng Cache L1d thứ 2 (Cấu trúc SoA Streams HTTP/2)\\[2pt]
    \normalfont\tiny \color{blue!80!black} \texttt{states[0..15]}: 16B $\bullet$ \texttt{window\_sizes[0..15]}: 32B $\bullet$ \texttt{flags[0..15]}: 16B
};

\node[draw=gray!50, fill=gray!10, dashed, minimum width=6.4cm, minimum height=0.55cm, font=\tiny, align=center] (coldcl) at (3.5, 0.4) {
    Phân vùng Cold Memory riêng biệt (Cấu hình TLS, Thống kê, Hostname)
};

\node[font=\scriptsize, align=center, text=teal!90!black] at (3.5, -0.4) {Bố cục nhớ phẳng, liên tục $\to$ Tận dụng 100\% băng thông Cache Line};
\end{scope}

\end{tikzpicture}
\caption{So sánh cấu trúc bộ nhớ: (a) Mô hình Hướng đối tượng truyền thống với mạng lưới con trỏ phân tán gây trượt cache nghiêm trọng; (b) Mô hình Định hướng Dữ liệu (\emph{Data-Oriented Design}) của $\mu$WebZockets gom nhóm dữ liệu nóng vào các dòng Cache Line 64-byte và tổ chức luồng theo Structure of Arrays (SoA).}
\label{fig:memory-layout}
\end{figure}

\paragraph{Kết quả dự kiến:}
\begin{itemize}
    \item Hoàn thiện module \path{src/core/pool.zig} quản lý Slab kết nối tĩnh và Freelist thời gian $O(1)$.
    \item Hoàn thiện cấu trúc SoA trong \path{src/http2/connection.zig} cho phép quản lý 100 luồng stream đồng thời trên mỗi kết nối HTTP/2 mà không tốn bất kỳ thao tác cấp phát heap nào trong thời gian chạy.
    \item Bảng thống kê kích thước bộ nhớ vật lý cho từng kết nối, chứng minh đạt chuẩn căn chỉnh 64 bytes (\texttt{align(64)}), loại bỏ 100\% hiện tượng False Sharing giữa các lõi CPU.
\end{itemize}

\subsubsection{Nội dung 2: Xây dựng cơ chế cô lập giao thức có giới hạn (Bounded Protocol Confinement) và hỗ trợ đa giao thức hiện đại.}

\paragraph{Phương pháp thực hiện:}
\begin{itemize}
    \item \textit{Nguyên tắc Cô lập Có giới hạn (\emph{Bounded Confinement}):} Quy định trần giới hạn tối đa cho mọi kích thước giao thức: \texttt{max\_request\_line\_size} ($1\text{ KB}$), \texttt{max\_header\_size} ($4\text{ KB}$), \texttt{max\_body\_size} ($2\text{ KB}\text{--}16\text{ KB}$), \texttt{max\_h2\_header\_block} ($4\text{ KB}$). Khi client gửi dữ liệu vượt quá ngưỡng quy định, parser ngay lập tức ngắt kết nối và phản hồi mã lỗi tương ứng ($414\text{ URI Too Long}$, $431\text{ Request Header Fields Too Large}$, $413\text{ Payload Too Large}$) thay vì tự động mở rộng bộ đệm gây cạn kiệt bộ nhớ hệ thống (\emph{Out-Of-Memory}).
    \item \textit{Kỹ thuật Mượn lát cắt không sao chép (\emph{Zero-Copy Borrowed Slices}):} Toàn bộ các đối tượng yêu cầu (\texttt{Request}) và phản hồi (\texttt{Response}) chỉ lưu trữ các con trỏ lát cắt (\texttt{[]const u8}) trỏ trực tiếp vào bộ đệm nhận của kết nối. Nhờ tính năng an toàn biên (\emph{spatial memory safety}) của Zig, không có dữ liệu nào bị sao chép thừa qua lại giữa tầng mạng và tầng ứng dụng.
    \item \textit{Bộ định tuyến Radix-Tree không cấp phát (\emph{Zero-Allocation Radix Router}):} Cây định tuyến URL được xây dựng tại thời điểm khởi động máy chủ hoặc biên dịch (\emph{comptime}). Khi khớp đường dẫn (\emph{route matching}), router chỉ trích xuất các tham số đường dẫn dưới dạng lát cắt bộ nhớ mượn (ví dụ: \path{/json/:count} mượn trực tiếp slice chuỗi số từ URL).
    \item \textit{Ngăn xếp Đa giao thức hoàn chỉnh:}
    \begin{itemize}
        \item \textbf{HTTP/1.1 \& RFC 10008 (HTTP QUERY):} Hỗ trợ đầy đủ pipelining, chunked transfer encoding, và phương thức an toàn có body \texttt{QUERY}~\cite{rfc10008}.
        \item \textbf{HTTP/2 \& HPACK:} Thực thi máy trạng thái kiểm tra preface, frame header, stream dependency, dynamic table giới hạn theo RFC 7541~\cite{rfc7541} và xử lý luồng mở rộng \texttt{CONNECT} (RFC 8441)~\cite{rfc8441}.
        \item \textbf{WebSocket \& Per-Message Deflate:} Tích hợp thư viện \texttt{zslay} do nhóm nghiên cứu phát triển cho phân tích khung WebSocket (RFC 6455)~\cite{rfc6455} và thư viện nén siêu tối ưu \texttt{libdeflate}~\cite{libdeflate} theo RFC 7692~\cite{rfc7692} với cơ chế đàm phán \texttt{no\_context\_takeover} để giải phóng bộ nhớ nén ngay sau mỗi khung thông điệp.
        \item \textbf{HTTP/3, QUIC \& WebTransport:} Tích hợp thư viện \texttt{lsquic}~\cite{lsquic} qua tầng FFI an toàn của Zig; hỗ trợ giải mã gói tin QUIC, quản lý luồng UDP bất đồng bộ, mở rộng Datagram không tin cậy (RFC 9221)~\cite{rfc9221} và bản nháp WebTransport-16~\cite{webtransport}.
        \item \textbf{Bảo mật TLS 1.3 \& ALPN:} Tích hợp BoringSSL~\cite{boringssl}, sử dụng cặp memory BIO kích thước cố định; cố tình tắt tính năng TLS 0-RTT nhằm ngăn chặn triệt để tấn công Replay Attack và bảo vệ biên quản lý bộ nhớ.
    \end{itemize}
\end{itemize}

\paragraph{Kết quả dự kiến:}
\begin{itemize}
    \item Hệ thống Router Radix-tree hoàn chỉnh có độ phức tạp tìm kiếm $O(K)$ với $K$ là độ dài URL, hoàn toàn không cấp phát heap trong quá trình định tuyến.
    \item Bộ kiểm thử Autobahn Testsuite ghi nhận đạt chuẩn $514$ ca \texttt{OK} và $3$ ca \texttt{INFORMATIONAL} trên $517$ ca kiểm thử, khẳng định độ tin cậy và sự tuân thủ chuẩn tuyệt đối của ngăn xếp WebSocket.
    \item Vượt qua 100\% các ca kiểm thử bảo mật biên: chống tràn số, chặn frame bất hợp lệ và từ chối các chuỗi tấn công DoS HTTP/2 CONTINUATION flood.
\end{itemize}

\subsubsection{Nội dung 3: Tối ưu hóa đường đi nóng (Hot Path) với vi lệnh SIMD và mô hình Multi-Worker Shared-Nothing.}

\paragraph{Phương pháp thực hiện:}
\begin{itemize}
    \item \textit{Gia tốc phân tích cú pháp tiêu đề bằng SIMD:} Trong giao thức HTTP/1.1, việc quét tìm ký tự kết thúc dòng \texttt{\textbackslash r\textbackslash n} (CRLF) và dấu hai chấm phân cách tiêu đề chiếm tới 30\% thời gian CPU của parser. Đề tài áp dụng kỹ thuật vector hóa SIMD (SWAR 64-bit trên kiến trúc cơ bản, AVX2 256-bit và ARM NEON 128-bit trên phần cứng hỗ trợ). Thuật toán tải cùng lúc 32 bytes dữ liệu vào thanh ghi vector và thực hiện phép so sánh song song trong một chu kỳ CPU duy nhất:
    \begin{equation}
        \mathbf{V}_{\text{match}} = \text{\_mm256\_cmpeq\_epi8}(\mathbf{V}_{\text{data}}, \mathbf{V}_{\text{target}})
    \end{equation}
    Sau đó sử dụng lệnh \texttt{\_mm256\_movemask\_epi8} kết hợp với lệnh đếm số bit 0 đầu tiên (\texttt{@ctz} trong Zig) để định vị tức thời vị trí của dấu phân cách mà không cần bất kỳ vòng lặp tuần tự nào.
    \item \textit{Thuật toán che mặt nạ WebSocket (SIMD Masking):} Theo RFC 6455, mọi frame từ client gửi lên server đều bị xáo trộn bởi một khóa mask 4-byte. Thay vì thực hiện phép XOR từng byte đơn lẻ, thuật toán nhân bản khóa 4-byte thành vector 32-byte và áp dụng lệnh \texttt{\_mm256\_xor\_si256}, cho phép xử lý giải mã payload với tốc độ chạm trần băng thông của thanh ghi CPU ($> 25\text{ GB/s}$).
    \item \textit{Mô hình Cụm Worker không chia sẻ (\emph{Shared-Nothing Cluster Architecture}):}
    Để tận dụng triệt để năng lực phần cứng đa lõi mà không vướng phải hiện tượng khóa nghẽn (\emph{lock contention}), hệ thống triển khai mô hình Multi-Worker:
    \begin{itemize}
        \item Số lượng worker được cấu hình tương ứng với số lõi CPU vật lý sẵn có ($4, 8, 16, 32, 64$ workers).
        \item Mỗi worker được ghim cố định (\emph{Core Affinity}) vào một lõi vi xử lý riêng biệt.
        \item Các worker cùng lắng nghe trên một cổng mạng chung thông qua cờ hệ điều hành \texttt{SO\_REUSEPORT}, cho phép nhân Linux tự động phân phối các kết nối TCP/UDP mới vào hàng đợi của từng lõi một cách cân bằng.
        \item Mỗi worker sở hữu một vòng lặp \texttt{libxev} riêng biệt với vòng đệm hoàn tất \texttt{io\_uring}, hoàn toàn không chia sẻ trạng thái hay sử dụng mutex/spinlock với các worker khác.
    \end{itemize}
\end{itemize}

\paragraph{Kết quả dự kiến:}
\begin{itemize}
    \item Tốc độ quét phân tích tiêu đề HTTP đạt trên $8.5\text{ GB/s}$ trên vi xử lý x86\_64 hỗ trợ AVX2, nhanh gấp 4 lần so với phương pháp lặp byte tuần tự.
    \item Tốc độ unmasking WebSocket đạt xấp xỉ tốc độ đọc bộ nhớ vật lý, giảm chi phí xử lý trên mỗi thông điệp xuống dưới $12\text{ ns}$.
    \item Hệ thống worker mở rộng tuyến tính gần như hoàn hảo ($> 95\%$ linear scaling efficiency) từ 1 lõi lên 8 lõi, 16 lõi và 32 lõi CPU.
\end{itemize}

\subsubsection{Nội dung 4: Hiện thực hóa bài toán trên bộ đo chuẩn quốc tế (bench.zig.bd, HttpArena) và đánh giá độ trễ đuôi bằng eBPF.}

\paragraph{Phương pháp thực hiện:}
\begin{itemize}
    \item \textit{Đo kiểm hiệu năng so sánh độc lập trên \url{https://bench.zig.bd/}~\cite{bench_zig_bd}:}
    Tiến hành đối sánh thông lượng (Requests/giây) và độ trễ của $\mu$WebZockets với toàn bộ các thư viện máy chủ mạng viết bằng Zig hiện có trên thị trường: \texttt{dusty+zio}, \texttt{httpz}, \texttt{zap}, \texttt{zzz} và máy chủ mặc định của thư viện chuẩn \texttt{std.http}. Thực nghiệm được tiến hành trên hai khối lượng công việc chuẩn:
    \begin{enumerate}
        \item \textbf{Workload Plaintext:} Đo năng lực I/O thuần túy và hiệu quả của bộ phân tích cú pháp HTTP framing.
        \item \textbf{Workload JSON:} Đo năng lực xử lý nghiệp vụ, bao gồm giải mã tham số URL, ánh xạ dữ liệu, tuần tự hóa JSON và nén dữ liệu động với Gzip/Deflate.
    \end{enumerate}
    \item \textit{Thẩm định qua 15 kịch bản tải nặng của HttpArena (PR \#1536)~\cite{httparena_pr1536}:}
    Triển khai $\mu$WebZockets trong môi trường container chuẩn hóa của HttpArena, thực thi bằng công cụ phát tải hạt mịn sử dụng \texttt{io\_uring} (\texttt{gcannon}) với số lượng kết nối đồng thời từ $256$, $512$, $1024$ đến $1536$ kết nối. Đo đạc các chỉ số: thông lượng đỉnh, tỷ lệ sử dụng CPU, dung lượng RSS RAM, và phân vị độ trễ (p50, p90, p99, p99.9).
    \item \textit{Đo kiểm hạt mịn độ trễ đuôi bằng Linux eBPF và Hardware Counters:}
    Để loại trừ hoàn toàn độ nhiễu đo kiểm từ phía client, nhóm nghiên cứu thiết lập các điểm hook eBPF trực tiếp trong nhân Linux~\cite{ebpf}:
    \begin{itemize}
        \item Tracepoint tại \texttt{sock:sock\_recvmsg} (thời điểm gói tin mạng vừa được đưa vào tầng socket).
        \item Uprobe tại điểm hoàn tất callback trong user-space của \texttt{libxev}.
        \item Tracepoint tại \texttt{sock:sock\_sendmsg} (thời điểm dữ liệu phản hồi được đẩy xuống card mạng).
        \item Sử dụng hàm \texttt{bpf\_ktime\_get\_ns()} để tính toán chính xác tuyệt đối thời gian xử lý thực sự (\emph{Service Time}) bên trong ứng dụng, ghi nhận vào BPF Ring Buffer mà không gây ảnh hưởng đến hiệu năng của hệ thống.
        \item Thu thập các sự kiện phần cứng bằng Linux \texttt{perf stat}: \texttt{L1-dcache-load-misses}, \texttt{L2-load-misses}, \texttt{branch-misses}, \texttt{instructions-per-cycle} (IPC).
    \end{itemize}
\end{itemize}

\begin{figure}[!hbt]
\centering
\begin{tikzpicture}[
    scale=0.92, transform shape,
    font=\footnotesize
]
% Axis
\draw[->, thick] (0,0) -- (12,0) node[below right, font=\small] {\textbf{Khung tải / Thư viện máy chủ}};
\draw[->, thick] (0,0) -- (0,5.5) node[above, font=\small] {\textbf{Thông lượng (ngàn req/s)}};

% Grid lines
\foreach \y/\val in {1/40, 2/80, 3/120, 4/160, 5/200} {
    \draw[gray!25, dashed] (0,\y) -- (11.5,\y);
    \node[left, font=\scriptsize] at (0,\y) {\val k};
}

% Plaintext Bars (Blue)
% uWebZockets: 170.01k -> 4.25
% dusty+zio: 165.20k -> 4.13
% httpz: 124.50k -> 3.11
% zap: 105.80k -> 2.64
% std.http: 82.40k -> 2.06
% zzz: 45.20k -> 1.13

\draw[fill=blue!70, draw=blue!90] (0.8,0) rectangle (1.4, 4.25) node[above, font=\tiny\bfseries] {170.0k};
\draw[fill=teal!60, draw=teal!80] (1.45,0) rectangle (2.05, 3.96) node[above, font=\tiny\bfseries] {158.4k};

\draw[fill=blue!70, draw=blue!90] (2.7,0) rectangle (3.3, 4.13) node[above, font=\tiny] {165.2k};
\draw[fill=teal!60, draw=teal!80] (3.35,0) rectangle (3.95, 3.80) node[above, font=\tiny] {152.1k};

\draw[fill=blue!70, draw=blue!90] (4.6,0) rectangle (5.2, 3.11) node[above, font=\tiny] {124.5k};
\draw[fill=teal!60, draw=teal!80] (5.25,0) rectangle (5.85, 2.81) node[above, font=\tiny] {112.3k};

\draw[fill=blue!70, draw=blue!90] (6.5,0) rectangle (7.1, 2.64) node[above, font=\tiny] {105.8k};
\draw[fill=teal!60, draw=teal!80] (7.15,0) rectangle (7.75, 2.35) node[above, font=\tiny] {94.2k};

\draw[fill=blue!70, draw=blue!90] (8.4,0) rectangle (9.0, 2.06) node[above, font=\tiny] {82.4k};
\draw[fill=teal!60, draw=teal!80] (9.05,0) rectangle (9.65, 1.77) node[above, font=\tiny] {71.0k};

\draw[fill=blue!70, draw=blue!90] (10.3,0) rectangle (10.9, 1.13) node[above, font=\tiny] {45.2k};
\draw[fill=teal!60, draw=teal!80] (10.95,0) rectangle (11.55, 0.97) node[above, font=\tiny] {38.9k};

% Labels under bars
\node[below=2pt, align=center, font=\scriptsize\bfseries] at (1.42, 0) {$\mu$WebZockets\\(Đề xuất)};
\node[below=2pt, align=center, font=\scriptsize] at (3.32, 0) {dusty+zio};
\node[below=2pt, align=center, font=\scriptsize] at (5.22, 0) {httpz};
\node[below=2pt, align=center, font=\scriptsize] at (7.12, 0) {zap};
\node[below=2pt, align=center, font=\scriptsize] at (9.02, 0) {std.http};
\node[below=2pt, align=center, font=\scriptsize] at (10.92, 0) {zzz};

% Legend
\node[draw=gray!60, fill=white, rounded corners=2pt, inner sep=4pt] at (8.5, 4.8) {
    \begin{tikzpicture}
        \draw[fill=blue!70, draw=blue!90] (0,0) rectangle (0.35,0.25);
        \node[right, font=\scriptsize] at (0.4, 0.12) {Plaintext Throughput (RPS)};
        \draw[fill=teal!60, draw=teal!80] (4.5,0) rectangle (4.85,0.25);
        \node[right, font=\scriptsize] at (4.9, 0.12) {JSON Throughput (RPS)};
    \end{tikzpicture}
};

\end{tikzpicture}
\caption{Đối sánh thông lượng xử lý yêu cầu (Requests/giây) giữa các máy chủ mạng viết bằng Zig theo số liệu công bố độc lập tại bộ chuẩn \textit{Zig Web Servers Benchmark} (\texttt{bench.zig.bd})~\cite{bench_zig_bd}. $\mu$WebZockets duy trì ưu thế dẫn đầu ở cả hai kịch bản tải Plaintext ($170{,}011\text{ req/s}$) và JSON ($158{,}420\text{ req/s}$).}
\label{fig:bench-throughput}
\end{figure}

\begin{figure}[!hbt]
\centering
\begin{tikzpicture}[
    scale=0.92, transform shape,
    font=\footnotesize
]
% Axis
\draw[->, thick] (0,0) -- (11.5,0) node[below right, font=\small] {\textbf{Số lượng kết nối đồng thời ($C$)}};
\draw[->, thick] (0,0) -- (0,5.0) node[above, font=\small] {\textbf{Độ trễ đuôi P99 ($\mu$s)}};

% Grid lines
\foreach \y/\val in {1/200, 2/400, 3/600, 4/800, 5/1000} {
    \draw[gray!25, dashed] (0,\y) -- (11,\y);
    \node[left, font=\scriptsize] at (0,\y) {\val};
}

% X-axis ticks
\foreach \x/\val in {1.8/64, 4.0/256, 6.2/512, 8.4/1024, 10.6/1536} {
    \draw (\x, 0.08) -- (\x, -0.08);
    \node[below, font=\scriptsize] at (\x, 0) {\val};
}

% Traditional AoS / Dynamic Heap line (Red)
% 64: 310us (1.55) -> 256: 480us (2.4) -> 512: 690us (3.45) -> 1024: 920us (4.6) -> 1536: 1250us (saturated > 5.0)
\draw[thick, draw=red!80, line width=1.4pt, mark=square*] plot coordinates {
    (1.8, 1.55)
    (4.0, 2.40)
    (6.2, 3.45)
    (8.4, 4.60)
    (10.6, 5.0)
};

% uWebZockets DoD Slab & SoA (Teal/Blue)
% 64: 180us (0.90) -> 256: 220us (1.10) -> 512: 290us (1.45) -> 1024: 380us (1.90) -> 1536: 450us (2.25)
\draw[thick, draw=teal!90!black, line width=1.4pt, mark=*] plot coordinates {
    (1.8, 0.90)
    (4.0, 1.10)
    (6.2, 1.45)
    (8.4, 1.90)
    (10.6, 2.25)
};

% Legend
\node[draw=gray!60, fill=white, rounded corners=2pt, inner sep=4pt] at (4.8, 4.4) {
    \begin{tikzpicture}
        \draw[thick, draw=teal!90!black, line width=1.2pt, mark=*] (0,0.12) -- (0.6,0.12);
        \node[right, font=\scriptsize] at (0.7, 0.12) {$\mu$WebZockets (Zero-Alloc Slab + SoA)};
        \draw[thick, draw=red!80, line width=1.2pt, mark=square*] (5.5,0.12) -- (6.1,0.12);
        \node[right, font=\scriptsize] at (6.2, 0.12) {Hệ thống AoS đối chứng (Heap Pointer-Chasing)};
    \end{tikzpicture}
};

\node[font=\scriptsize, align=left, fill=teal!10, draw=teal!70, rounded corners=2pt] at (8.5, 1.2) {
    \textbf{Đặc tính phẳng ổn định:}\\
    Nhờ loại bỏ GC và cấp phát heap,\\
    độ trễ P99 chỉ tăng nhẹ từ $180\mu\text{s}$ lên\\
    $450\mu\text{s}$ ngay cả khi tải tăng gấp 24 lần.
};

\end{tikzpicture}
\caption{Biểu đồ độ trễ đuôi phân vị P99 theo số lượng kết nối đồng thời từ $64$ đến $1536$ kết nối (dữ liệu trích xuất từ kịch bản đo kiểm eBPF socket hook và bộ chuẩn \textit{HttpArena} PR \#1536~\cite{httparena_pr1536}). $\mu$WebZockets duy trì đường cong độ trễ phẳng, ngăn chặn hiện tượng bùng nổ độ trễ ở tải nặng.}
\label{fig:tail-latency}
\end{figure}

\begin{table}[!ht]
\centering
\small
\caption{Tổng hợp ma trận năng lực và kết quả kiểm thử đa giao thức trong \textit{HttpArena} (PR \#1536) và \textit{Autobahn Testsuite}}
\label{tab:protocol-matrix}
\begin{tabularx}{\textwidth}{@{}lXll@{}}
\toprule
\textbf{Giao thức / Thành phần} & \textbf{Tiêu chuẩn IETF / Hiện thực hóa} & \textbf{Cơ chế quản lý bộ nhớ} & \textbf{Kết quả thực nghiệm} \\
\midrule
HTTP/1.1 Core & RFC 9112, RFC 10008 (QUERY) & Bounded Slab, Zero-Alloc & $170{,}011\text{ req/s}$ (Pass 4/4) \\
HTTP/2 Multiplexing & RFC 9113, HPACK RFC 7541 & SoA Stream Slabs, Huffman Comptime & $100\%$ Pipelined OK \\
WebSocket RFC 6455 & zslay parser, SIMD Masking & Fixed Frame Ring Buffer & 514 OK / 3 Info (Autobahn) \\
Per-Message Deflate & RFC 7692 via libdeflate & \texttt{no\_context\_takeover} Bounded & Pass echo-ws-pipeline \\
HTTP/3 \& QUIC & RFC 9114, RFC 9000 via lsquic & Bounded Packet/Stream Pool & Pass baseline-h3 \\
WebTransport & draft-ietf-webtrans-http3-16 & Session Slab, Bounded Datagrams & Khung kiểm thử hoàn tất \\
Bảo mật TLS 1.3 & RFC 8446, ALPN via BoringSSL & Fixed Memory BIO, 0-RTT Rejected & Pass json-tls, json-h2c \\
\bottomrule
\end{tabularx}
\end{table}

\paragraph{Kết quả dự kiến:}
\begin{itemize}
    \item Bộ dữ liệu thực nghiệm đo đạc độc lập toàn diện được công bố công khai, chứng minh $\mu$WebZockets giảm hơn $65\%$ tỷ lệ trượt cache L1d/L2 so với máy chủ đối chứng viết theo mô hình đối tượng thông thường.
    \item Độ trễ đuôi phân vị P99 và P99.9 duy trì ở mức dưới $1\text{ ms}$ trên các mức tải cực lớn lên đến $1536$ kết nối đồng thời và $500{,}000$ requests trong bài test \texttt{latency-500k-8cpu} của HttpArena.
    \item Báo cáo phân tích chuyên sâu về tương quan giữa số chu kỳ CPU trên mỗi request với tỷ lệ tiêu thụ thông lượng mạng thực tế.
\end{itemize}

\subsection{B2.3 Kế hoạch nghiên cứu}

Kế hoạch nghiên cứu được phân chia thành 5 giai đoạn nối tiếp chặt chẽ trong vòng 12 tháng:

\begin{longtable}[c]{@{} >{\raggedright\arraybackslash}p{(\columnwidth - 6\tabcolsep) * \real{0.12}} >{\raggedright\arraybackslash}p{(\columnwidth - 6\tabcolsep) * \real{0.24}} >{\raggedright\arraybackslash}p{(\columnwidth - 6\tabcolsep) * \real{0.38}} >{\raggedright\arraybackslash}p{(\columnwidth - 6\tabcolsep) * \real{0.26}}@{}}
\toprule()
\textbf{Giai đoạn} & \textbf{Thời gian} & \textbf{Nội dung công việc} & \textbf{Sản phẩm dự kiến} \\
\midrule()
\endhead
Giai đoạn 1 & Tháng 1 -- 2 & Khảo sát vi kiến trúc CPU x86/ARM, cơ chế cache L1/L2, kernel \texttt{io\_uring} và mã nguồn Zig 0.16. & Báo cáo tổng quan lý thuyết; Bản thiết kế cấu trúc Slab và SoA. \\
Giai đoạn 2 & Tháng 3 -- 5 & Hiện thực hóa core engine: Connection Slab, Freelist $O(1)$, parser SIMD, Radix Router và tích hợp \texttt{libxev}. & Mã nguồn \path{src/core/}; Thư viện module Zig chạy thử nghiệm. \\
Giai đoạn 3 & Tháng 6 -- 8 & Hiện thực hóa đa giao thức: HTTP/2 (SoA), WebSocket (zslay, libdeflate), QUIC/H3 (lsquic), BoringSSL TLS 1.3. & Bộ parser đa giao thức; Báo cáo Autobahn Testsuite (514 OK / 517). \\
Giai đoạn 4 & Tháng 9 -- 10 & Triển khai đo kiểm chuẩn hóa trên \url{https://bench.zig.bd/}, HttpArena (15 kịch bản) và profiling eBPF. & Bộ dữ liệu đo đạc chi tiết; Biểu đồ phân vị latency và cache miss. \\
Giai đoạn 5 & Tháng 11 -- 12 & Đánh giá, tối ưu tinh chỉnh hệ thống, viết báo cáo tổng kết đề tài và bài báo khoa học công bố. & 01 bài báo khoa học; Báo cáo toàn văn nghiệm thu đề tài NCKHSV. \\
\bottomrule()
\end{longtable}

\subsection{B3. Kết quả dự kiến}

\paragraph{1. Sản phẩm khoa học và công bố:}
\begin{itemize}
    \item \textbf{Bài báo khoa học:} 01 bài báo khoa học được chấp nhận đăng tại Kỷ yếu Hội nghị Khoa học Quốc tế có phản biện (như IEEE RIVF, NICS hoặc Tạp chí Phát triển Khoa học \& Công nghệ -- ĐHQG-HCM), với nội dung phân tích chuyên sâu về kiến trúc Data-Oriented Design và kỹ thuật tối ưu hóa bộ nhớ đệm cho máy chủ mạng bất đồng bộ.
    \item \textbf{Báo cáo nghiên cứu toàn văn:} Báo cáo thuyết minh và tổng kết khoa học đề tài NCKHSV năm 2026, trình bày tường minh cơ sở lý thuyết, mô hình kiến trúc, mã nguồn hiện thực hóa và phương pháp luận đo kiểm.
\end{itemize}

\paragraph{2. Sản phẩm ứng dụng và mã nguồn công nghệ:}
\begin{itemize}
    \item \textbf{Thư viện mã nguồn mở $\mu$WebZockets:} Kho mã nguồn mở chất lượng cao trên GitHub, đóng gói thành Zig Package hoàn chỉnh, có tài liệu hướng dẫn API đầy đủ (\texttt{Typed Request/Response}) cho cộng đồng lập trình viên Zig quốc tế và sinh viên các khóa sau học tập.
    \item \textbf{Bộ công cụ đo kiểm và dữ liệu Benchmark chuẩn hóa:} Tích hợp chính thức trong hệ sinh thái \textit{Zig Web Servers Benchmark} (\texttt{bench.zig.bd}) và \textit{HttpArena} (PR \#1536), kèm theo các kịch bản kiểm thử độ trễ hạt mịn dựa trên Linux eBPF có khả năng tái lập độc lập.
\end{itemize}

\paragraph{3. Sản phẩm đào tạo và tác động thực tiễn:}
\begin{itemize}
    \item Nâng cao năng lực nghiên cứu hệ thống cấp thấp, lập trình mạng hiệu năng cao, hiểu biết về vi kiến trúc CPU và công nghệ kernel Linux hiện đại cho sinh viên thực hiện đề tài thuộc Khoa Mạng máy tính và Truyền thông dữ liệu, Trường ĐH Công nghệ Thông tin.
    \item Tạo nguồn học liệu mẫu chất lượng cao cho sinh viên các môn học chuyên ngành như: Hệ điều hành nâng cao, Lập trình mạng, Tối ưu hóa hiệu năng hệ thống phân tán.
\end{itemize}

\nocite{*}
\printbibliography[heading=bibintoc, title={B4. Tài liệu tham khảo}]

\vspace{1cm}
\centering
\begin{longtable}[c]{@{} >{\centering\arraybackslash}p{(\columnwidth - 2\tabcolsep) * \real{0.6}} >{\centering\arraybackslash}p{(\columnwidth - 2\tabcolsep) * \real{0.4}}@{}}
% \toprule()
\begin{minipage}[b]{\linewidth}\centering \emph{Ngày \ldots tháng \ldots năm 2026}\\ \textbf{Giảng viên hướng dẫn}\\ (Ký và ghi rõ họ tên) \end{minipage} &
\begin{minipage}[b]{\linewidth}\centering \emph{Ngày \ldots tháng \ldots năm 2026}\\ \textbf{Chủ nhiệm đề tài}\\ (Ký và ghi rõ họ tên) \end{minipage} \\
\end{longtable}

\end{document}
